\section{Код в Python для решения уравнения осцилляции нейтрино в среде}
\label{app: code-python-vis}

Для численного решения уравнения \eqref{eq:15} использовался код на python. 


Ниже представлен код сценария использованного для получения данных по квазипериодам \(\Re \psi_1(\xi)\), \(\Im\psi_1(\xi)\) в зависимости от \(\xi\), а также данных по количеству переходов через ось абсцисс \(\Re\psi_{2, 3}\) и \(\Im\psi_{2, 3}\) на интервале каждого из ранее найденных квазипериодов \(\Re\psi_1\) и \(\Im\psi_1\) соответственно. 

Его также можно найти в репозитории \url{https://github. com/Fest-Fut/Kursovay}.

\inputminted[linenos,breaklines,fontsize=\scriptsize]{python}{../run.py}

%\textattachfile{stat-prop_v2.py}{Код в виде прикреплённых данных (только для
%	электронной версии документа).}
